\section{Limitations and Threats to Validity}

The first limitation is dataset access. The project is designed around the Dartmouth campus WLAN trace, but the raw movement archive was not directly available through legacy unauthenticated URLs during implementation. The included evaluation therefore uses a deterministic trace-compatible fallback. This validates the pipeline but does not replace evaluation on the raw Dartmouth archive.

The second limitation is synthetic content demand. Zipf popularity and spatial locality are research-backed assumptions, but real application demand may include burstiness, temporal drift, content lifetimes, and user-specific preferences. Future work should rerun the pipeline with real content logs if available.

The third limitation is scale. The final local experiment uses 20 nodes and 48 objects. Larger graphs would require faster candidate evaluation, batched shortest-path computation, and possibly graph neural encoders. The current implementation favors clarity and reproducibility over large-scale optimization.

The fourth limitation is selector dependence on validation perturbations. Adaptive TACC assumes that held-out perturbation samples resemble the test regime. If topology changes in a qualitatively different way, the selector may choose a suboptimal policy. This can be addressed with online monitoring, confidence intervals, and periodic retraining.

Finally, the DQN uses explicit graph features rather than a full GNN. This keeps the model lightweight, but a GNN may generalize better across graph sizes and node permutations. The present method should therefore be seen as a strong baseline and research scaffold for future graph-learning extensions.

